\begin{PSbox}[unbreakable]{obj}{اهداف یادگیری}

این ماژول پایانی، تمام ماژول‌های پیشین را در کاربردهای واقعی مهندسی برق یکپارچه می‌کند. دانشجویان قادر خواهند بود:

\begin{enumerate}
\def\labelenumi{\arabic{enumi}.}
\tightlist
\item
  مدل‌سازی احتمالاتی را در سیستم‌های مخابراتی به کار ببرند
\item
  مسایل پردازش سیگنال را به‌صورت آماری تحلیل کنند
\item
  تغییرات قطعات الکترونیکی را مدل کنند
\item
  قابلیت اطمینان سیستم‌های قدرت را ارزیابی کنند
\item
  عدم‌قطعیت اندازه‌گیری در سیستم‌های کنترل را مدیریت کنند
\item
  احتمال/آمار را به مبانی یادگیری ماشین مرتبط کنند
\end{enumerate}

\end{PSbox}

\PSsection{16.1}{مخابرات}

\PSsubsection{\PSnum{16.1.1} برآورد \lr{\mbox{BER}} و بازه‌های اطمینان}

\textbf{مسیله:} برآورد نرخ خطای بیت یک لینک مخابراتی.

\textbf{مدل:} هر بیت یک آزمایش برنولی با احتمال خطای \lr{$p$} است (همان \lr{\mbox{BER}}).

\begin{itemize}
\tightlist
\item
  \lr{$n$} بیت ارسال‌شده، \lr{$k$} خطا مشاهده‌شده
\item
  \lr{$\hat{p} = \frac{k}{n}$} (برآوردگر \lr{\mbox{MLE}} برای \lr{\mbox{BER}})
\item
  بازه اطمینان 95\%: \lr{$\hat{p} \pm 1.96\sqrt{\frac{\hat{p}(1-\hat{p})}{n}}$}
\end{itemize}

\PSsubsection{\PSnum{16.1.2} مدل‌سازی نویز \lr{\mbox{AWGN}}}

نویز گاوسی سفید جمع‌شونده \lr{\mbox{(AWGN)}} مدل نویز بنیادی است:

\begin{itemize}
\tightlist
\item
  \textbf{توزیع:} \lr{$N\left(0,\, \frac{N_{0}}{2}\right)$} در هر بُعد
\item
  \textbf{چرا گاوسی؟} \lr{\mbox{CLT}} \lr{—} نویز حرارتی جمع سهم‌های بسیاری از الکترون‌ها است
\item
  \textbf{مشخصه‌سازی:} کاملاً با چگالی طیفی توان \lr{$N_{0}$} یا واریانس \lr{$\sigma^{2} = N_{0}B$} توصیف می‌شود (\lr{$B$} = پهنای‌باند)
\end{itemize}

\PSsubsection{\PSnum{16.1.3} کانال محوشدگی رایلی}

\textbf{مدل:} ضریب کانال \lr{$h = h_{I} + j \cdot h_{Q}$} که در آن \lr{$h_{I},\, h_{Q} \sim N(0,\, \sigma^{2})$} و مستقل هستند.

\begin{itemize}
\tightlist
\item
  پوش \lr{$\lvert h\rvert \sim \mathrm{Rayleigh}(\sigma)$}
\item
  توان \lr{$\lvert h\rvert^{2} \sim \mathrm{Exponential}\left(\frac{1}{2\sigma^{2}}\right)$}
\item
  \textbf{احتمال قطع:} \lr{$P(\lvert h\rvert^{2} < \gamma_{\mathrm{th}}) = 1 - \exp \left(-\frac{\gamma_{\mathrm{th}}}{2\sigma^{2}}\right)$}
\end{itemize}

\PSsubsection{\PSnum{16.1.4} محوشدگی رایسی}

هنگامی که مؤلفه خط دید \lr{\mbox{(LOS)}} وجود دارد:

\begin{itemize}
\tightlist
\item
  \lr{$h = h_{\mathrm{LOS}} + h_{\mathrm{scatter}}$}
\item
  \lr{$\lvert h\rvert \sim \mathrm{Rician}(\nu,\, \sigma)$} که در آن \lr{$K = \frac{\nu^{2}}{2\sigma^{2}}$} ضریب \lr{$K$} است
\end{itemize}

ضریب \lr{\mbox{$K$: $K = 0$}} (رایلی) \lr{$\to K = \infty$} (بدون محوشدگی، \lr{\mbox{AWGN}})

\PSsubsection{\PSnum{16.1.5} آماره‌های \lr{\mbox{SNR}}}

\textbf{\lr{\mbox{SNR}} لحظه‌ای:} \lr{$\gamma = \lvert h\rvert^{2} \times \frac{E_{s}}{N_{0}}$}

برای رایلی: \lr{$\gamma \sim \mathrm{Exponential}(\bar{\gamma})$} که در آن \lr{\mbox{$\bar{\gamma}$ = SNR}} میانگین

\textbf{میانگین \lr{\mbox{BER}} در محوشدگی:}\PSdisplay{\bar{P}_e = \int_0^\infty P_e(\gamma) \cdot f_\gamma(\gamma) \, d\gamma}

برای \lr{\mbox{BPSK}} در کانال رایلی: \lr{$\bar{P}_{e} = \frac{1}{2}\left(1 - \sqrt{\frac{\bar{\gamma}}{1+\bar{\gamma}}}\right)$}

\PSsection{16.2}{پردازش سیگنال}

\PSsubsection{\PSnum{16.2.1} مشخصه‌سازی نویز اندازه‌گیری}

\textbf{مسیله:} مشخصه‌سازی نویز در سیگنال اندازه‌گیری‌شده \lr{$x(t) = s(t) + n(t)$}.

\textbf{رویکرد:}

\begin{enumerate}
\def\labelenumi{\arabic{enumi}.}
\tightlist
\item
  نمونه‌های فقط‌نویز (در غیاب سیگنال) را جمع کنید
\item
  آماره‌های نویز را برآورد کنید: \lr{$\hat{\mu}_{n}$}، \lr{$\hat{\sigma}^{2}_{n}$}
\item
  گاوسی‌بودن را آزمون کنید (هیستوگرام، نمودار \lr{\mbox{QQ}})
\item
  \lr{\mbox{SNR}} را محاسبه کنید: \lr{$\mathrm{SNR} = \frac{P_{\mathrm{signal}}}{\hat{\sigma}^{2}_{n}}$}
\end{enumerate}

\PSsubsection{\PSnum{16.2.2} آشکارسازی مبتنی بر همبستگی}

\textbf{مسیله:} آشکارسازی یک سیگنال معلوم در نویز.

\textbf{روش:} سیگنال دریافتی را با الگوی معلوم همبسته کنید. تحت \lr{$H_{0}$} (فقط نویز)، خروجی همبسته‌ساز توزیع معلومی دارد \lr{\mbox{\ensuremath{\to}}} آستانه را تعیین کنید.

\PSsection{16.3}{الکترونیک}

\PSsubsection{\PSnum{16.3.1} مدل‌سازی رواداری قطعه}

\textbf{مسیله:} مقاومت‌ها مقدار نامی \lr{$R_{0}$} با رواداری ساخت دارند.

\textbf{مدل:} \lr{$R \sim N(R_{0},\, \sigma_{R}^{2})$} که در آن \lr{$\sigma_{R}$} به رده رواداری بستگی دارد.

برای رواداری \lr{$\pm 5\%$} (\PSnum{99.7\%} در محدوده): \lr{$3\sigma = 0.05R_{0} \to \sigma = 0.0167R_{0}$}

\PSsubsection{\PSnum{16.3.2} تحلیل مدار با مونت‌کارلو}

\textbf{مسیله:} یک تقسیم‌کننده ولتاژ از \lr{$R_{1}$} و \lr{$R_{2}$} استفاده می‌کند. توزیع ولتاژ خروجی چیست؟

\lr{$\displaystyle V_{\mathrm{out}} = V_{\mathrm{in}} \times \frac{R_{2}}{R_{1} + R_{2}}$}

با تصادفی‌بودن \lr{$R_{1}$} و \lr{\mbox{$R_{2}$: $V_{\mathrm{out}}$}} یک متغیر تصادفی است!

\PSsubsection{\PSnum{16.3.3} انتشار عدم‌قطعیت اندازه‌گیری}

اگر \lr{$Z = f(X,\, Y)$} و \lr{$X$}، \lr{$Y$} عدم‌قطعیت‌های کوچک داشته باشند:

\PSdisplay{\sigma_Z^2 \approx \left(\frac{\partial f}{\partial X}\right)^2 \sigma_X^2 + \left(\frac{\partial f}{\partial Y}\right)^2 \sigma_Y^2}

\textbf{مثال:} توان \lr{$P = \frac{V^{2}}{R}$. $\sigma_{P}^{2} \approx \frac{2\,\mathrm{V}}{R}^{2}\sigma_{V}^{2} + \frac{V^{2}}{R^{2}}^{2}\sigma_{R}^{2}$}

\PSsection{16.4}{سیستم‌های قدرت}

\PSsubsection{\PSnum{16.4.1} عدم‌قطعیت بار}

\textbf{مسیله:} بار اوج نامعین است. مدل: \lr{$L \sim N(\mu_{L},\, \sigma_{L}^{2})$}

پرسش طراحی: چه ظرفیت \lr{$C$} تضمین می‌کند \lr{$P(\text{\rl{بار}} > \text{\rl{ظرفیت}}) < 0.01$}؟

\lr{$C = \mu_{L} + 2.326\sigma_{L}$} (برای احتمال تخطی 1\%)

\PSsubsection{\PSnum{16.4.2} قابلیت اطمینان سیستم}

\textbf{سیستم سری:} همه قطعات باید کار کنند.

\lr{$R_{\mathrm{sys}} = R_{1} \times R_{2} \times \ldots \times R_{n}$} (برای قطعات مستقل)

\textbf{سیستم موازی:} حداقل یکی باید کار کند.

\lr{$R_{\mathrm{sys}} = 1 - (1-R_{1})(1-R_{2})\ldots (1-R_{n})$}

\PSsection{16.5}{سیستم‌های کنترل}

\PSsubsection{\PSnum{16.5.1} اثرهای نویز حسگر}

\textbf{مسیله:} یک سیستم کنترل از قرایت‌های نویزی حسگر استفاده می‌کند.

اندازه‌گیری: \lr{$y(k) = x(k) + n(k)$} که در آن \lr{$n \sim N(0,\, \sigma^{2}_{n})$}

کنترل‌کننده روی \lr{$y(k)$} عمل می‌کند \lr{\mbox{\ensuremath{\to}}} خطای ردیابی شامل سهم نویز می‌شود.

\textbf{رویکردهای حل:}

\begin{itemize}
\tightlist
\item
  میانگین‌گیری (فیلتر پایین‌گذر): نویز را به اندازه \lr{$\sqrt{n}$} کاهش می‌دهد
\item
  فیلتر کالمن: تلفیق بهینه پیش‌بینی مدل و اندازه‌گیری
\end{itemize}

\PSsubsection{\PSnum{16.5.2} عدم‌قطعیت پارامتر}

مدل سیستم: \lr{$G(s) = \frac{K}{\tau s + 1}$} که در آن \lr{$K$} و \lr{$\tau$} نامعین هستند.

رویکرد مونت‌کارلو: \lr{$K$} و \lr{$\tau$} را از توزیع‌هایشان نمونه‌گیری کنید، پاسخ سیستم را بارها شبیه‌سازی کنید و آماره‌های خروجی را محاسبه کنید.

\PSsection{16.6}{ارتباط با یادگیری ماشین / هوش مصنوعی}

\PSsubsection{\PSnum{16.6.1} نمونه‌گیری و جمع‌آوری داده}

\begin{itemize}
\tightlist
\item
  داده آموزش = نمونه‌ای از جامعه
\item
  داده آزمون = نمونه مستقل برای اعتبارسنجی
\item
  فرض \lr{\mbox{i.i.d}}. زیربنای بیشتر نظریه یادگیری ماشین است
\end{itemize}

\PSsubsection{\PSnum{16.6.2} برآورد به‌عنوان بنیاد یادگیری ماشین}

\begin{itemize}
\tightlist
\item
  رگرسیون خطی = \lr{\mbox{MLE}} تحت فرض نویز گاوسی
\item
  آموزش شبکه عصبی = \lr{\mbox{MLE}} (یا \lr{\mbox{MAP}}) برای پارامترهای مدل
\item
  کمینه‌سازی تابع هزینه = اصل بیشینه درست‌نمایی
\end{itemize}

\PSsubsection{\PSnum{16.6.3} توزیع‌های احتمال در یادگیری ماشین}

\begin{itemize}
\tightlist
\item
  مدل‌های مولد: یادگیری مستقیم \lr{$P(X)$}
\item
  مدل‌های مخلوط گاوسی: جمع وزنی گاوسی‌ها
\item
  خودرمزگذارهای تغییرپذیر: متغیرهای نهفته با پیشین‌های گاوسی
\end{itemize}

\PSsubsection{\PSnum{16.6.4} توازن اریبی-واریانس}

\PSdisplay{\PSmtext{Test Error} = \PSmtext{Bias}^2 + \PSmtext{Variance} + \PSmtext{Irreducible Noise}}

این دقیقاً همان تجزیه \lr{\mbox{MSE}} از نظریه برآورد است (ماژول 13)!

\begin{itemize}
\tightlist
\item
  اریبی زیاد: کم‌برازش (مدل بسیار ساده)
\item
  واریانس زیاد: بیش‌برازش (مدل بسیار پیچیده)
\item
  بهینه: توازن میان این دو
\end{itemize}

\PSsubsection{\PSnum{16.6.5} اطمینان و عدم‌قطعیت در پیش‌بینی‌ها}

\begin{itemize}
\tightlist
\item
  بازه‌های اطمینان برای پیش‌بینی‌های مدل
\item
  شبکه‌های عصبی بیزی: توزیع پسین روی وزن‌ها \lr{\mbox{\ensuremath{\to}}} عدم‌قطعیت پیش‌بینی
\item
  کالیبراسیون: احتمال‌های پیش‌بینی‌شده باید با فراوانی‌های مشاهده‌شده همخوانی داشته باشند
\end{itemize}

\PSsection{16.7}{نقاط یکپارچه‌سازی کلیدی}

{\def\LTcaptype{none} % do not increment counter
\begin{longtable}[c]{@{}cc@{}}
\toprule\noalign{}
\rowcolor{PSnavy}\PSth{ماژول} & \PSth{کاربرد در مهندسی برق} \\
\midrule\noalign{}
\endhead
\bottomrule\noalign{}
\endlastfoot
1. احتمال & احتمال‌های رویداد، تشخیص عیب \\
2. متغیرهای تصادفی & مدل‌سازی کمیت‌های نامعین \\
3. توزیع‌ها & انتخاب مدل‌های نویز/محوشدگی/خرابی \\
4. گشتاورها & توان سیگنال، توان نویز، \lr{\mbox{SNR}} \\
5. تبدیل‌ها & تبدیل به \lr{\mbox{dB}}، توان از ولتاژ \\
\lr{$6-7$}. مشترک/توابع & سیگنال+نویز، ترکیب‌های سیستم \\
8. شرطی & عملکرد در وضعیت‌های مختلف \\
9. همبستگی & کانال‌های \lr{\mbox{MIMO}}، آرایه‌های حسگر \\
10. چند متغیر تصادفی & پردازش برداری، شکل‌دهی پرتو \\
11. \lr{\mbox{CLT}} & چرا نویز گاوسی است، میانگین‌گیری \\
12. آمار & برآورد پارامترهای سیستم \\
13. برآورد & \lr{\mbox{MLE}} برای پارامترهای کانال/نویز \\
14. اطمینان & کمّی‌سازی عدم‌قطعیت اندازه‌گیری \\
15. فرض & کنترل کیفیت، مقایسه عملکرد \\
\end{longtable}
}

\emph{پایان ماژول‌های درس. ماژول‌های درس با کاربردهای مهندسی و جمع‌بندی مفاهیم کلیدی پایان می‌یابند.}
